\section{APU}

Emulace subsystému APU byla nejtěžší etapou ve vývoji emulátoru, jelikož ladit zvuk není tak
jednoduchý jako grafický výstup -- v případě PPU jakékoliv nesrovnatelnosti jsou snadno
pozorovatelné a ve většině času příčina chyby je přibližně jasná. V případě zvuku je to poněkud
komplikovanější. Proto většinu času při emulaci APU bylo stráveno laděním a opětovným přečtením
dostupné dokumentace.

Nejdříve byl implementován pulzní kanál, protože je nejjednodušší a nejpoužívanější v demonstračních
hrách. Postupně se implementovaly ostatní kanály. Desktopový backend využívá zvukový API knihovny
SDL3, zatímco backend RP2350 využívá vlastní zvukový ovladač, který bude popsán v sekci
\ref{pp:audio}.

\subsection{Vzorkování (sampling)}

Jak již bylo zmíněno v \ref{sub:downsampling}, je potřeba provést downsampling v případě jestli
zvuková jednotka generuje vzorky s příliš velkou frekvenci, což je případ APU v NESu.

Použijeme metodu decimace, která spočívá ve filtraci signálu pomocí dolní propustí a následném
výběru $N$-tého vzorku. Naše implementace ale zvolila metodu jednodušší decimaci -- metodu přímého
výběru každého $N$-tého vzorku bez dolní propustí. Vzhledem k nelineární povaze a nízkému rozlišení
původního zvukového čipu APU je výsledný rozdíl v kvalitě zvuku pro běžného uživatele téměř
nepostřehnutelný.

Pro realizaci výběru $N$-tého vzorku, je uchováván akumulátor času mezi vzorky. Při každém vydaném
vzorku se k akumulátoru přičte čas. Problémem ale je, že čas mezi jednotlivými vzorky nemusí být
celočíselný -- závisí na taktovací frekvencí emulovaného CPU $f_{CPU}$ (pro systém NES platí
$f_{CPU} \approx \SI{1.79}{\mega\hertz}$) a zvolenem sample ratu $f_s$, což lze vyjadřit vztahem:

$$
step = \frac{f_{CPU}}{f_s}
$$

Pro RP2350 backend platí $f_s = \SI{31.25}{\kilo\hertz}$ (desktop backend používá standardní frekvenci vzorkovaní
44,1 kHz). V případě backendu RP2350 vychází tento krok přibližně na 57,1786 cyklů.
Proto použití zaokrouhlování pro akumulátor času by vedlo k akumulaci chyby v časování, což by se
projevilo slyšitelným kolísáním výšky tónu nebo praskáním vlivem nesprávné fáze signálu.

Aby bylo možné dosáhnout optimální přesnosti bez použití výpočetně náročné aritmetiky s plovoucí čárkou,
využívá implementace \textbf{aritmetiku s pevnou řádovou čárkou} (\textit{fixed-point}) ve formátu
16.16. V tomto formátu je 32bitové slovo rozděleno na dvě poloviny:

\begin{itemize}
    \item \textbf{celočíselnou část} -- představuje celočíselný počet cyklů.
    \item \textbf{zlomkovou část} -- uchovává desetinnou část s přesností na $\frac{1}{65536}$.
\end{itemize}

\subsection{Krokování a generování vzorků}

Funkce \texttt{apu\_run\_until} (viz výpis \ref{lst:apu_run_until}) aktualizuje stav APU a posouvá
ho po cyklové ose dokud nedožene stav CPU (parametr \texttt{target\_cycle}).

\begin{listing}[htbp]
\begin{minted}{c}
void apu_run_until(Apu *self, CycleCounter target_cycle) {
    CycleCounter cycles_to_run = target_cycles - self->cycles;

    while (cycles_to_run >= 2) {
        /* Aktualizace všech kanálů (trojúhelníkový kanál se
         * aktualizuje 2krát, jelikož je taktován stejnou frekvenci
         * jako CPU) */
        update_triangle(&self->triangle);
        update_triangle(&self->triangle);
        update_pulse(&self->p1);
        update_pulse(&self->p2);
        update_noise(&self->noise);
        update_frame_counter(self);
        /* Downsampling: vzorek se generuje jen po překročení
         * určitého času (po poslední generaci) */
        self->sample_acc += (2 << 16); /* +2 cyklů CPU */
        if (self->sample_acc >= self->sample_step) {
            self->sample_acc -= self->sample_step;
            output_sample(self);
        }

        cycles_to_run -= 2;
        self->cycles += 2;
    }
}
\end{minted}
\caption{Aktualizace stavu APU.}
\label{lst:apu_run_until}
\end{listing}

Při každém kroku APU se aktualizuji vnitřní stav a výstup každého kanálů. Následně se ověřuje, zda
je na čase vygenerovat nový vzorek -- jakmile akumulátor \texttt{sample\_acc} dosáhne prahové
hodnoty \texttt{sample\_step}, je vygenerován výstupní vzorek a z akumulátoru je odečtena délka
kroku, čimž je zajištěna fázová spojitost signálu bez kumulativní chyby v časování.

Ačkoliv je APU taktováno stejnou frekvencí jako CPU, většina jeho komponentů je aktualizována pouze
každý druhý cyklus CPU. V rámci implementace je proto výhodné zpracovávat APU po blocích dvou CPU
cyklů, přičemž komponenty běžící na plné frekvenci (trojúhelníkový kanál) jsou aktualizovány dvakrát
v rámci jednoho kroku.

Aby bylo možné minimalizovat drahé operace dělení a aritmetiku plovoucí čárky, vzorky se generují pomocí LUT
tabulky (viz \cite{nesdev-apu-lut}). Tyto LUT tabulky se inicializuji při prvotní inicializaci APU (výpis
\ref{lst:apu_sample_output}).

\begin{listing}[htbp]
\begin{minted}{c}
void apu_init(Apu *self, Bus *bus) {
    /* Zkráceno... */
    for (int i = 0; i < 31; ++i) {
        pulse_lut[i] = (95.52 / (8128.0 / (double) i + 100.0)) *
                                    APU_SAMPLE_MAX_VALUE;
    for (int i = 0; i < 203; ++i) {
        tnd_lut[i] = (163.67 / (24329.0 / (double) i + 100)) *
                                    APU_SAMPLE_MAX_VALUE;
}

void output_sample(Apu *self) {
    /* Zkráceno (ověřování jestli jednotlivé kanály jsou zapnuté
     * nebo vypnuté)... */
    ApuSample pulse_smp = pulse_lut[p1_out + p2_out];
    ApuSample tnd_smp = tnd_lut[3 * tri_out + 2 * noise_out + dmc_out];

    self->sample_cb(self, pulse_smp + tnd_smp);
}
\end{minted}
\caption{Generace vzorků pomocí LUT tabulky, která má předvypočítané hodnoty pomocí lineární
aproximace popsané v \cite{nesdev}.}
\label{lst:apu_sample_output}
\end{listing}
